\chapter{Complexe onderdelen van EMAV}
\label{ch:projectcomplexiteit}

Hoewel de overstap naar een modern webframework (.NET Razor Pages) veel van de technologische beperkingen van het oude Windows Forms-project wegneemt, blijft EMAV-web een applicatie met een aanzienlijke cognitieve last voor ontwikkelaars. Deze complexiteit is niet beperkt tot de architecturale keuzes, maar is diep geworteld in de aard van het domein en de technische keuzes die nodig zijn om aan de wetenschappelijke eisen te voldoen.

\section{Domeincomplexiteit: De kloof tussen wetenschap en implementatie}
Het meest prominente obstakel bij het begrijpen en beheren van de codebase is niet zozeer dat de programmeurs de logica niet zouden vatten — zij zijn immers verantwoordelijk voor de vertaling van de wiskundige modellen naar werkende software — maar eerder de enorme drempel die de C\#-implementatie opwerpt voor de onderzoekers zelf. In tegenstelling tot klassieke administratieve applicaties vormt de kern van EMAV een web van onderling afhankelijke wetenschappelijke formules die gebaseerd zijn op jarenlang landbouwkundig onderzoek.

Hoewel de code voor een ontwikkelaar technisch helder en logisch gestructureerd kan zijn, fungeert ze voor een onderzoeker vaak als een black box. De wetenschappers, die de eigenlijke domeinexperts zijn, herkennen hun eigen wiskundige modellen vaak niet onmiddellijk terug in de abstracties van een objectgeoriënteerde taal. Dit zorgt voor een fundamenteel probleem in de levenscyclus van het project: wanneer een onderzoeker een formule wil finetunen of een nieuwe parameter wil testen, is hij volledig afhankelijk van de IT-afdeling. Het maken van zelfs de kleinste aanpassingen vereist een cyclus van overleg, implementatie en validatie door de programmeur, wat experimentatie en snelle iteratie ernstig belemmert. De complexiteit zit dus niet in de onbegrijpelijkheid van de code voor IT, maar in het gebrek aan autonomie voor de wetenschapper binnen een puur technische omgeving.

\section{Hybride Databeheer en Persistentiestrategieën}
{\sloppy Een ander complex onderdeel is de manier waarop EMAV omgaat met data. Het systeem maakt gebruik van een hybride architectuur waarbij verschillende database-technologieën naast elkaar worden gebruikt. Metadata over gebruikers, projecten en overzichten worden opgeslagen in een centrale SQL Server-database via Entity Framework Core. De ruwe inputgegevens voor berekeningen bevinden zich echter vaak in SQLite-bestanden die per project of per onderzoeksgroep worden beheerd.\par}

Deze scheiding is noodzakelijk om onderzoekers de nodige flexibiliteit te bieden om met eigen datasets te werken, maar het introduceert aanzienlijke complexiteit in de datatoegangslaag. Ontwikkelaars moeten navigeren tussen verschillende database providers, rekening houden met inconsistente schema's en zorgen voor data-integriteit over de grenzen van de verschillende databases heen. Het beheer van deze persistentiestromen vereist een diepgaand inzicht in zowel de infrastructuurlaag als de domeinregels.

\section{De \textquotedbl Legacy Bridge\textquotedbl \@ en Technische Schuld}
Hoewel dat EMAV-web als een nieuw project is gestart, is het onvermijdelijk dat er logica uit de oude applicatie is overgenomen. Legacy-systemen bevatten vaak bedrijfskritische kennis die niet zomaar mag verdwijnen; volgens \textcite{comellaDorda2000survey} zijn zulke systemen vaak tegelijk te belangrijk om weg te gooien en te kwetsbaar om zonder meer te wijzigen. Het hergebruiken van complexe, gevalideerde rekenmodules uit de Windows Forms-tijd is daarom essentieel om de wetenschappelijke consistentie te garanderen. Dit creëert echter een Legacy Bridge: een zone in de code waar moderne C\#-conventies en verouderde programmeerpatronen met elkaar in conflict komen.

Het adapteren van deze oude logica naar een moderne, asynchrone Clean Architecture vereist veel plumbing code. Deze aanpak sluit aan bij incrementele modernisering: \textcite{seacord2001legacyModernization} beschrijven dat modernisering vaak betekent dat legacy-componenten en vernieuwde componenten tijdelijk naast elkaar moeten functioneren, wat bijkomende integratieproblemen introduceert. Ook wrapping en componentization worden in de literatuur genoemd als technieken om bestaande legacy-functionaliteit via moderne interfaces beschikbaar te maken \autocite{comellaDorda2000survey}. Ontwikkelaars worden geconfronteerd met methoden die niet voldoen aan de huidige standaarden voor testbaarheid of leesbaarheid, maar die vanwege hun kritieke rol in de berekeningen niet zomaar herschreven kunnen worden. Deze voorzichtigheid is herkenbaar in het werk van \textcite{feathers2004legacyCode}, waar legacy-code benaderd wordt als code waarvan het bestaande gedrag eerst beschermd moet worden voordat ze veilig kan worden aangepast. Deze mix van oud en nieuw verhoogt de drempel voor nieuwe ontwikkelaars om de volledige flow van de applicatie te begrijpen.

\section{Container-Orchestratie en Omgevingsbeheer}
Tot slot vormt de integratie met Docker en Jupyter Notebooks een technische uitdaging. Zoals uitgebreid besproken in Hoofdstuk \ref{ch:poc}, start de applicatie op runtime containers op voor de onderzoekers. Dit vereist dat de webapplicatie directe controle heeft over de Docker-daemon. Docker waarschuwt in de eigen documentatie dat de daemon normaal met root-rechten werkt en dat toepassingen die via een API containers aanmaken daarom bijzonder voorzichtig moeten omgaan met invoer en parameters \autocite{dockerEngineSecurity}. Hierdoor is container-orchestratie in EMAV niet louter een implementatiedetail, maar een beveiligingskritisch onderdeel van de architectuur.

Daarnaast introduceert het netwerkmodel extra complexiteit. Containers zijn standaard geïsoleerd, maar het publiceren van poorten doorbreekt die isolatie door forwarding-regels naar de host te creëren \autocite{dockerPublishingPorts}. Wanneer bij port publishing geen expliciet IP-adres wordt opgegeven, kan Docker de poort bovendien standaard op alle interfaces publiceren \autocite{dockerRunReference}. Voor een Jupyter-omgeving die via tokens toegankelijk is, betekent dit dat poortbindingen, netwerkbereik en tokenbeheer samen correct geconfigureerd moeten worden.

De noodzaak om DLL's dynamisch in te laden in een Linux-container, poortbindingen te beheren, resource-limieten in te stellen en tokens te extraheren uit logs, voegt een laag van systeem-engineering toe aan de applicatie. Ook resourcebeheer is hier relevant: Docker vermeldt dat containers zonder expliciete limieten standaard zoveel CPU en geheugen kunnen gebruiken als de kernel scheduler toelaat \autocite{dockerResourceConstraints}. Dit soort logica bevindt zich op het raakvlak tussen softwareontwikkeling en systeembeheer, wat het debugging-proces bemoeilijkt: een fout kan immers zowel in de C\#-code, de Docker-configuratie als in de Python-omgeving van de container liggen.

{\sloppy Om deze uiteenlopende vormen van complexiteit — variërend van de wetenschappelijke detaillering tot de technische gelaagdheid van de data — beheersbaar te houden, is een ad-hoc programmeerstijl ontoereikend. In complexe softwareprojecten is het net belangrijk om verantwoordelijkheden expliciet te scheiden en afhankelijkheden beheerst te organiseren; dit sluit aan bij de principes van Clean Architecture, waarbij de businesslogica beschermd wordt tegen details zoals frameworks, infrastructuur en externe technologieën \autocite{martin2017clean}. In plaats van losse, situatiegebonden oplossingen steunt de architectuur van EMAV-web daarom op een doordacht frame van software-ontwerppatronen (design patterns). Design patterns beschrijven volgens \textcite{gof1994} herbruikbare oplossingen voor terugkerende ontwerpproblemen in objectgeoriënteerde software. Deze patronen dienen dus niet louter als theoretische exercitie, maar vormen een direct antwoord op de hierboven geschetste uitdagingen.\par}

Zo is het Repository patroon een noodzakelijk instrument om de hybride datastructuur te temmen, terwijl CQRS en het Mediator patroon de cognitieve last voor ontwikkelaars verlagen door verantwoordelijkheden strikt te scheiden. Tegelijkertijd biedt het Strategy patroon de nodige wetenschappelijke flexibiliteit binnen de rekenkern, en fungeert Dependency Injection als de lijm die al deze losgekoppelde componenten op een testbare manier samenbrengt. Deze keuze past binnen het bredere idee dat patronen en architectuurprincipes niet op zichzelf staan, maar gebruikt worden om domeinlogica onderhoudbaar, testbaar en technologisch minder afhankelijk te maken \autocite{millett2015patterns,seemann2019di}. In de volgende subsecties worden de belangrijkste patronen in detail besproken, waarbij telkens de link wordt gelegd tussen het theoretische concept en de concrete meerwaarde voor de robuustheid en onderhoudbaarheid van EMAV.

\section{Repository Pattern}
{\sloppy Het repository patroon fungeert als een essentiële abstractielaag tussen de domeinlogica en de diverse gegevensbronnen waarover EMAV beschikt. Volgens \textcite{fowlerRepository} gedraagt een repository zich als een in-memory verzameling van domeinobjecten, waarbij de technische details van de onderliggende persistentiemechanismen volledig worden geabstraheerd. In de context van EMAV is dit patroon cruciaal omdat het project opereert op een hybride datastructuur die over de jaren heen is geëvolueerd. Enerzijds is er een SQLite-database voor de opslag van ruwe inputgegevens en rekendata, en anderzijds een SQL Server-database voor de webgerelateerde metadata en gebruikersbeheer.\par}

Binnen de codebase wordt dit gerealiseerd via een gelaagde hiërarchie van interfaces en implementaties. 
In het \texttt{ILVO.EMAV.Shared}-project bevindt zich de basisdefinitie \texttt{IBaseRepository<T>}, die algemene CRUD-operaties (Create, Read, Update, Delete) definieert. 
Specifieke projecten breiden dit uit, zoals \texttt{IWebBaseRepository<TEntity>} in de core-laag, waarbij gebruik wordt gemaakt van generieke constraints (where TEntity \: class, IWebBaseEntity). 
De \texttt{ILVO.EMAV.Core}-laag communiceert uitsluitend met deze interfaces. 
Hierdoor blijft de kern van de applicatie volledig onwetend over de specifieke implementatie details, of die nu gebruikmaakt van Entity Framework Core in \texttt{ILVO.EMAV.Infrastructure.Web} of gespecialiseerde logic in de input-laag.

Deze abstractie verhoogt de onderhoudbaarheid aanzienlijk: indien een database-schema wijzigt of indien men besluit over te stappen naar een andere opslagtechnologie, hoeven de wijzigingen enkel in de infrastructuurlaag te worden doorgevoerd. 
De businesslogica blijft onaangetast. 
Bovendien is dit patroon de hoeksteen voor kwalitatieve unit testing. 
Door repositories te mocken, kan de businesslogica worden getoetst zonder afhankelijkheid van een actieve databaseverbinding. 
Dit wijst op een architecturale keuze die gericht is op stabiliteit op lange termijn en eenvoudige verifieerbaarheid van de code \autocite{millett2015patterns}.

\section{Command Query Responsibility Segregation (CQRS)}
{\sloppy CQRS is een geavanceerd architectuurpatroon waarbij de verantwoordelijkheid voor het opvragen van gegevens (Queries) strikt wordt gescheiden van de verantwoordelijkheid voor het wijzigen van de systeemtoestand (Commands). Volgens \textcite{fowlerCQRS} biedt deze scheiding de mogelijkheid om elk pad onafhankelijk te structureren, te optimaliseren en te schalen. In een complex systeem zoals EMAV, waar berekeningen uren kunnen duren en metadata-opvragingen milliseconden moeten duren, is deze scheiding bittere noodzaak.\par}

Binnen EMAV wordt CQRS geïmplementeerd met de \texttt{LiteBus}-bibliotheek \autocite{litebus}. 
Dit resulteert in een zogenaamde Vertical Slice Architecture binnen de core-projecten \autocite{jimmyBogardVerticalSlice}. 
Een actie zoals het verwijderen van een berekening wordt niet afgehandeld door een generieke CalculationService, maar door een specifieke \texttt{DeleteCalculationCommand} en de bijbehorende \texttt{DeleteCalculationCommandHandler}. 
Dit commando bevat enkel de minimale data (zoals een \texttt{Guid Id}) en de handler voert de noodzakelijke stappen uit : het ophalen van de metadata via een repository, het verwijderen van fysieke bestanden via een \texttt{IFileManager} en het finaal verwijderen van de database-records.

Aan de query-zijde worden specifieke modellen (DTO's) geretourneerd die geoptimaliseerd zijn voor de weergave in de UI, in plaats van de volledige domeinentiteiten te lekken. 
Deze aanpak voorkomt het ontstaan van Gouden Klassen die alle logica van een bepaald domein bevatten. 
In plaats van dat de logica gedistribieerd over honderden kleine, gefocuste klassen die elk één specifieke taak hebben. 
Dit verlaagt de cognitieve last voor ontwikkelaars die aan het project werken aanzienlijk.

\section{Mediator Pattern}
Onlosmakelijk verbonden met de CQRS-aanpak is het mediator patroon. 
Dit patroon fungeert als een centrale verkeersleider die de communicatie tussen verschillende objecten coördineert, waardoor directe koppelingen tussen componenten worden vermeden \autocite{gof1994}. 
In de moderne architectuur van EMAV is the \texttt{IMediator} (geleverd door LiteBus) het enige aanspreekpunt voor de presentatielaag (\texttt{ILVO.EMAV.Web}).

Een Razor Page in de UI-laag hoeft geen enkel inzicht te hebben in welke services, repositories of validators nodig zijn om een bepaalde actie uit te voeren. 
De pagina stuurt simpelweg een commando of query naar de mediator. 
De mediator zorgt er vervolgens voor dat dit bericht door de juiste pipelines passeert. 
Zo worden commando's automatisch eerst gevalideerd door \texttt{IValidator}-implementaties voordat ze de handler bereiken.

Dit zorgt voor een extreme vorm van ontkoppeling : de presentatielaag heeft geen enkele directe referentie naar de infrastructuur- of datalaag. 
Dit verhoogt de flexibiliteit van het systeem enorm; handlers kunnen worden herschreven, verplaatst naar andere microservices of gedecoreerd met extra logging zonder dat er ook maar één regel code in de UI-laag hoeft te wijzigen. 
Dit patroon is een directe indicatie van de architecturale volwassenheid van het project en draagt bij aan de robuustheid tegen technologische veranderingen.

\section{Strategy Pattern}
Het strategie patroon is een gedragspatroon dat toelaat om een algoritme te selecteren op runtime en dit algoritme in te kapselen in een uitwisselbaar object \autocite{gof1994}. 
In de rekenkern van EMAV is dit patroon van fundamenteel belang om de wetenschappelijke variabiliteit van de emissiemodellen te ondersteunen. 
De kerninterface \texttt{ICalculator<Tin, Tout>} vormt de basis for een uitgebreide familie van rekenstrategieën.

Deze architectuur is essentieel omdat EMAV berekeningen uitvoert voor diverse domeinen, zoals enterische emissies, stalemissies en mestverwerking. 
Elke berekening heeft haar eigen set formules, inputparameters en validatieregels. 
Door elke berekening als een afzonderlijke strategie (bijvoorbeeld \texttt{BerekenEmissieLandbouwbodem}) te implementeren, blijft de overkoepelende orchestratielogica clean en overzichtelijk. 
De orchestrator roept simpelweg de \texttt{DoCalc}-methode aan op de geïnjecteerde strategie.

Dit bevordert het Open-Closed Principle (onderdeel van SOLID): de rekenkern staat open voor uitbreiding (men kan eenvoudig een nieuwe reken module toevoegen door de interface te implementeren), maar is gesloten voor aanpassingen (de bestaande orchestratie hoeft niet gewijzigd te worden). 
Voor een project waar wetenschappers regelmatig nieuwe formules aanleveren, is dit een cruciale eigenschap die voorkomt dat het project na verloop van tijd vastloopt in haar eigen complexiteit.

\section{Factory Pattern}
Het factory patroon wordt aangewend om de creatie van objecten te abstraheren van de client-code. 
Dit is met name nuttig wanneer het exacte type van het aan te maken object afhankelijk is van configuratie of gebruikersinput op runtime. 
In EMAV wordt dit patroon prominent toegepast in de \texttt{IExporterFactory}.

Het systeem ondersteunt een breed scala aan exportmogelijkheden, waaronder CSV voor verdere analyse in statistische software, Excel voor rapportage en JSON voor integratie met andere systemen. {\sloppy De \texttt{ExporterFactory} neemt een \texttt{ExportType} als parameter en retourneert de juiste implementatie van de \texttt{IExporter}-interface. Hierdoor hoeft de aanroepende code niet te weten hoe een specifieke exporter geconfigureerd moet worden.\par}

Deze centralisatie van creatielogica voorkomt codeduplicatie en maakt het systeem zeer modulair. 
Indien er in de toekomst ondersteuning voor een nieuw formaat moet worden toegevoegd, hoeft enkel een nieuwe klasse te worden geschreven en geregistreerd in de factory. 
De rest van de applicatie blijft volledig ongewijzigd. 
Dit illustreert de visie van het project op uitbreidbaarheid en het vermijden van harde koppelingen tussen componenten \autocite{gof1994}.

\section{Facade Pattern}
Gezien de enorme interne complexiteit van de EMAV-rekenkern, waarbij tientallen subsystemen gelijktijdig actief zijn, is het facade patroon onontbeerlijk. 
Een facade biedt een vereenvoudigde interface naar een complex subsysteem, waardoor de interne complexiteit wordt afgeschermd voor de gebruiker \autocite{gof1994}. 
In EMAV fungeren de belangrijkste controllers en orchestratie-services als facades.

Wanneer een gebruiker op Start Berekening klikt, lijkt dit voor de UI een enkele actie. 
Achter de facade vindt echter een complexe dans plaats: inputbestanden worden gevalideerd, de juiste rekenfactoren worden uit de database opgehaald, cache-mechanismen worden geraadpleegd, meerdere berekeningsstrategieën worden in de juiste volgorde aangeroepen en de resultaten worden finaal geaggregeerd en weggeschreven. 
De facade verbergt deze interacties en biedt een stabiele API.

Dit is een erkenning van de inherente complexiteit van het domein: in plaats van te proberen de complexiteit volledig weg te nemen, wordt ze beheersbaar gemaakt door ze te groeperen achter logische, gebruiksvriendelijke interfaces. Dit verbetert de onderhoudbaarheid omdat de interne werking van de rekenkern kan worden geoptimaliseerd zonder dat de UI-code hiervan hinder ondervindt.

\section{Dependency Injection en Inversion of Control}
{\sloppy Dependency Injection is de architecturale ruggengraat die alle voorgaande patronen in EMAV faciliteert en verbindt. In plaats van dat klassen hun eigen afhankelijkheden instantiëren, worden deze afhankelijkheden van buitenaf aangeboden. Volgens \textcite{seemann2019di} is DI de enige manier om een echt modulaire en testbare software-architectuur te bouwen.\par}

EMAV maakt optimaal gebruik van de ingebouwde DI-container van .NET. 
{\sloppy Elk deelproject (Core, Infrastructure, Web) bevat een eigen \nolinkurl{DependencyInjection.cs}-bestand waarin de interne services, handlers en repositories worden geregistreerd. Dit bevordert een hoge mate van inkapseling: de UI-laag hoeft enkel de DI-methode van de Core-laag aan te roepen om alle benodigde logica beschikbaar te maken.\par}

Dit mechanisme is cruciaal voor de flexibiliteit van het project. 
Het laat toe om op een transparante manier implementaties te wisselen op basis van de omgeving of voor testdoeleinden. 
DI is de motor achter de andere patronen; zonder een goede DI-structuur zouden patronen zoals Strategy en Repository leiden tot een onhoudbare hoeveelheid handmatige object-creatie doorheen de applicatie.

\section{Adapter Pattern}
Het adapter patroon maakt het mogelijk dat klassen met incompatibele interfaces toch naadloos kunnen samenwerken \autocite{gof1994}. 
In de context van EMAV is dit patroon een cruciaal instrument om de kloof tussen de moderne Clean Architecture en de waardevolle legacy-code overbruggen. 
Zoals eerder beschreven, bevat het project logica die overgenomen is uit een oudere Windows Forms-omgeving.

In plaats van deze oude code integraal te herschrijven, worden adapters ingezet. 
Deze adapters presenteren de oude functionaliteit aan de nieuwe applicatie via een interface die voldoet aan de nieuwe standaarden (bijvoorbeeld door gebruik te maken van \texttt{Task<T>} voor asynchroon werken). 
Hierdoor kan de nieuwe kernlogica modern en clean blijven, terwijl ze toch gebruikmaakt van de bewezen nauwkeurigheid van de bestaande rekenmodules. 
Dit toont een pragmatische en architecturaal verantwoorde aanpak van software-evolutie aan.

\section{Decorator Pattern en Result Pattern}
Het decorator patroon voegt dynamisch verantwoordelijkheden toe aan een object zonder de onderliggende klasse te wijzigen \autocite{gof1994}. 
In de EMAV-codebase wordt dit principe toegepast via de mediator pipeline. 
Elke actie die door het systeem stroomt, wordt gedecoreerd met extra functionaliteit zoals automatische validatie via \texttt{FluentValidation}. 
Hierdoor blijven de eigenlijke business-handlers puur en gefocust op hun specifieke taak.

Daarnaast wordt het Result Pattern (via de \texttt{ErrorOr} bibliotheek) consistent toegepast. 
In plaats van uitzonderingen (exceptions) te gebruiken voor control flow bij te verwachten fouten, retourneren methoden een object dat ofwel het resultaat, ofwel een lijst van fouten bevat. 
Dit patroon, zoals beschreven door \textcite{khorikovResultPattern}, zorgt voor een zeer voorspelbare en robuuste foutafhandeling die naadloos integreert met de Clean Architecture-filosofie en functionele programmeerprincipes.

\section{Clean Architecture}
Alle voorgaande patronen komen samen in de overkoepelende Clean Architecture-structuur van EMAV. 
Volgens \textcite{martin2017clean} is het hoofddoel van deze architectuur om de businesslogica volledig te isoleren van infrastructuur, frameworks en externe invloeden. 
In de projectstructuur van EMAV is deze filosofie strikt doorgevoerd:
\begin{itemize}
    \item Domain Layer (\nolinkurl{ILVO.EMAV.Data.*}): Bevat de pure entiteiten en domeinobjecten.
    \item Application Layer (\texttt{ILVO.EMAV.Core}): Bevat de businessregels, de mediator handlers en de definities van de interfaces.
    \item Infrastructure Layer (\nolinkurl{ILVO.EMAV.Infrastructure.*}): Bevat de technische implementaties, zoals database-connectiviteit en bestandssysteem-interacties.
    \item Presentation Layer (\texttt{ILVO.EMAV.Web}): De user interface (Razor Pages) die uitsluitend communiceert met de Core-laag.
\end{itemize}

Deze gelaagdheid zorgt ervoor dat EMAV een zeer hoge mate van technologische onafhankelijkheid geniet. De Clean Architecture is de ultieme uiting van de complexiteitsbeheersing binnen dit project: het biedt een duurzaam frame waarin alle andere design patterns optimaal tot hun recht komen om samen een kwalitatief superieur softwareproduct te vormen.
